\documentclass[11pt]{article}
\usepackage{amsmath,amsthm,amssymb}
\usepackage[margin=1in]{geometry}

\newtheorem{theorem}{Theorem}
\newtheorem{lemma}{Lemma}
\newtheorem{proposition}{Proposition}
\newtheorem{corollary}{Corollary}
\theoremstyle{definition}
\newtheorem{remark}{Remark}
\newtheorem{definition}{Definition}

\DeclareMathOperator{\Des}{Des}
\DeclareMathOperator{\maj}{maj}
\DeclareMathOperator{\comaj}{comaj}
\DeclareMathOperator{\SYT}{SYT}
\DeclareMathOperator{\ord}{ord}
\DeclareMathOperator{\core}{core}
\newcommand{\ZZ}{\mathbb{Z}}
\newcommand{\la}{\lambda}
\newcommand{\inner}[1]{\left\langle #1 \right\rangle}

\title{The graded $2$-core vanishing law:\\
$\displaystyle \ord_{q=-1} G_\lambda(q) = \big\lfloor |\,2\text{-}\core(\lambda)\,|/2 \big\rfloor$}
\author{Clio}
\date{2026-06-01}

\begin{document}
\maketitle

\begin{abstract}
For a partition $\la\vdash n$ let
$G_\la(q)=\sum_{T\in\SYT(\la)} q^{s(T)}$, where
$s(T)=\sum_{i\in\Des(T)}w_i$ with parity-twisted branch weights
$w_i=2i-1$ if $n-i$ is odd and $w_i=0$ otherwise. This is the graded generating
function of the branch-exponent (Theorem~B) descent statistic; it is generically
irreducible over $\mathbb{Q}$ and has \emph{no} product formula. I prove that its
order of vanishing at $q=-1$ is governed entirely by the $2$-core of $\la$:
\[
   \ord_{q=-1}G_\la(q)\;=\;c(\la)\;:=\;\Big\lfloor \tfrac12\,|2\text{-}\core(\la)|\Big\rfloor,
\]
and I give the first nonzero Taylor coefficient at $q=-1$ in closed form via the
$2$-quotient. The proof rests on a single closed identity expressing $G_\la(-q)$ as
a Schur--inner product,
\[
  G_\la(-q)=\Big\langle s_\la,\; h_1^{\,e}\!\!\prod_{k\,:\,n-k\ \mathrm{odd}}\!\!\big(h_2-e_2\,q^{\,2k-1}\big)\Big\rangle,
  \qquad e=n\bmod 2,
\]
after which the law follows from one structural fact of character theory: removing
dominoes preserves the $2$-core, so $\chi^\la(2^a1^b)=0$ whenever
$b<|2\text{-}\core(\la)|$. Specialising at $q=1$ recovers the previously established
fiber $G_\la(-1)=\chi^\la(w_0)$. All steps are verified computationally for every
partition of $n\le 12$.
\end{abstract}

\section{Statement}

Throughout $\la\vdash n$ and $\SYT(\la)$ is the set of standard Young tableaux of
shape $\la$. For $T\in\SYT(\la)$,
\[
  \Des(T)=\{\,i\in[n-1]: i{+}1\text{ is in a strictly lower row of }T\text{ than }i\,\},
  \qquad \comaj(T)=\sum_{i\in\Des(T)}(n-i).
\]
Fix the \emph{active set} and the branch weights
\[
  A=\{\,k\in[n-1]: n-k\text{ odd}\,\},\qquad
  w_k=(2k-1)\,[\,k\in A\,],\qquad
  s(T)=\sum_{i\in\Des(T)}w_i,
\]
and put $m=|A|=\lfloor n/2\rfloor$ and $e=n\bmod 2$ (so $A$ is the set of even
positions when $n$ is odd, the set of odd positions when $n$ is even). The object of
study is
\[
  G_\la(q)=\sum_{T\in\SYT(\la)}q^{s(T)}\ \in\ \ZZ_{\ge0}[q].
\]
Recall that every partition has a $2$-core (the unique partition with no removable
domino reachable by deleting dominoes) of size $c_0:=|2\text{-}\core(\la)|$ and a
$2$-quotient $(\la^{(0)},\la^{(1)})$; one has $c_0\equiv n\pmod 2$ and
$|\la^{(0)}|+|\la^{(1)}|=w:=(n-c_0)/2$ (the number of dominoes). Set
$c(\la)=\lfloor c_0/2\rfloor$, so that $c_0=2c+e$.

\begin{theorem}\label{thm:main}
For every $\la\vdash n$,
\[
  \ord_{q=-1}G_\la(q)\;=\;c(\la)\;=\;\Big\lfloor \tfrac12\,|2\text{-}\core(\la)|\Big\rfloor.
\]
Writing $G_\la(q)=\sum_{k\ge0}a_k\,(q+1)^k$, the first nonzero coefficient is
\[
  a_{c}\;=\;(-1)^{c}\,e_{c}\!\big(\{2k-1:k\in A\}\big)\cdot\Big(-\tfrac12\Big)^{c}\,
           \chi^\la\!\big(2^{\,m-c}\,1^{\,c_0}\big),
\]
where $e_c$ is the elementary symmetric polynomial of degree $c$. Equivalently, via
the $2$-quotient,
\[
  a_{c}=(-1)^{n(\la)}\,f^{\la^{(0)}}f^{\la^{(1)}}\binom{w}{|\la^{(0)}|}\,S(n,c),
  \qquad
  S(n,c)=\pm\,e_c\!\big(\{2k-1:k\in A\}\big)\,2^{-c}\,f^{\,2\text{-}\core(\la)},
\]
a scalar depending only on $(n,c)$; in particular $S(n,0)=1$ and
$S(n,1)=-\binom n2$.
\end{theorem}

Here $f^\mu=\#\SYT(\mu)$ and $n(\la)=\sum_i(i-1)\la_i$.

\section{The master identity}

Let $\Lambda$ be the ring of symmetric functions, $s_\la$ the Schur functions,
$h_r,e_r,p_r$ the complete, elementary and power-sum bases, and
$\inner{\,\cdot\,,\cdot\,}$ the Hall inner product, for which
$\inner{s_\la,h_\mu}=K_{\la\mu}$ (the Kostka number) and
$\inner{s_\la,p_\mu}=\chi^\la(\mu)$. For a subset $R\subseteq[n-1]$ with
$R=\{r_1<\dots<r_t\}$ let $\alpha(R)=(r_1,r_2-r_1,\dots,n-r_t)$ be the composition of
$n$ with descent set $R$. We use the classical fact
\begin{equation}\label{eq:kostka-des}
  K_{\la,\alpha(R)}=\inner{s_\la,h_{\alpha(R)}}=\#\{T\in\SYT(\la):\Des(T)\subseteq R\}
\end{equation}
(Stanley, \emph{EC2}, \S7.19; the right equality is standardisation of semistandard
tableaux of content $\alpha(R)$). Write $\beta_\la(S)=\#\{T:\Des(T)=S\}$, so that
$K_{\la,\alpha(R)}=\sum_{S\subseteq R}\beta_\la(S)$.

\begin{lemma}[Descent enumerator]\label{lem:descent}
For indeterminates $\eta=(\eta_1,\dots,\eta_{n-1})$,
\[
  \Psi_\la(\eta):=\sum_{T\in\SYT(\la)}\ \prod_{k\in\Des(T)}\eta_k
   \;=\;\sum_{R\subseteq[n-1]} K_{\la,\alpha(R)}\ \prod_{k\in R}\eta_k\prod_{k\notin R}(1-\eta_k).
\]
\end{lemma}

\begin{proof}
Substitute $K_{\la,\alpha(R)}=\sum_{S\subseteq R}\beta_\la(S)$ into the right-hand
side and exchange the order of summation:
\[
  \sum_{R}\Big(\sum_{S\subseteq R}\beta_\la(S)\Big)\prod_{k\in R}\eta_k\prod_{k\notin R}(1-\eta_k)
  =\sum_{S}\beta_\la(S)\sum_{R\supseteq S}\prod_{k\in R}\eta_k\prod_{k\notin R}(1-\eta_k).
\]
For fixed $S$, writing $R=S\sqcup S'$ with $S'\subseteq[n-1]\setminus S$,
\[
  \sum_{R\supseteq S}\prod_{k\in R}\eta_k\prod_{k\notin R}(1-\eta_k)
  =\prod_{k\in S}\eta_k\prod_{k\in[n-1]\setminus S}\big(\eta_k+(1-\eta_k)\big)
  =\prod_{k\in S}\eta_k .
\]
Hence the right-hand side equals $\sum_S\beta_\la(S)\prod_{k\in S}\eta_k=\Psi_\la(\eta)$.
\end{proof}

\begin{proposition}[Master identity]\label{prop:master}
As polynomials in $q$,
\[
  G_\la(-q)\;=\;\Big\langle s_\la,\; h_1^{\,e}\prod_{k\in A}\big(h_2-e_2\,q^{\,2k-1}\big)\Big\rangle .
\]
\end{proposition}

\begin{proof}
Set $\eta_k=(-1)^{n-k}q^{\,w_k}$ in Lemma~\ref{lem:descent}. Since
$(-1)^{n-k}=(-1)^{w_k}$ for every $k$ (indeed $w_k$ is odd exactly when $n-k$ is
odd), $\prod_{k\in\Des(T)}\eta_k=(-1)^{\comaj(T)}q^{s(T)}=(-1)^{s(T)}q^{s(T)}=(-q)^{s(T)}$,
so the left side of Lemma~\ref{lem:descent} is $\Psi_\la(\eta)=G_\la(-q)$.

For $k\notin A$ we have $n-k$ even and $w_k=0$, hence $\eta_k=1$ and the factor
$1-\eta_k$ vanishes; therefore only subsets $R\supseteq A^c$ contribute. Write
$R=A^c\sqcup U$ with $U\subseteq A$. For $k\in A$ we have $\eta_k=-q^{2k-1}$ and
$1-\eta_k=1+q^{2k-1}$, while for $k\in A^c$ the factor $\eta_k=1$ is harmless. Thus
\begin{equation}\label{eq:U-sum}
  G_\la(-q)=\sum_{U\subseteq A}K_{\la,\alpha(A^c\cup U)}\,
            \prod_{k\in U}\big(-q^{2k-1}\big)\prod_{k\in A\setminus U}\big(1+q^{2k-1}\big).
\end{equation}

\emph{Content of $\alpha(A^c\cup U)$.} The marks of $\alpha(A^c)$ are
$\{0\}\cup A^c\cup\{n\}$. When $n$ is even, $A^c=\{2,4,\dots,n-2\}$ and the gaps are
all equal to $2$, giving the multiset $(2^m)$; when $n$ is odd,
$A^c=\{1,3,\dots,n-2\}$ and the gaps are $(1,2^m)$. In either case the content is
$(2^m,1^e)$, and the $m$ parts equal to $2$ occupy the gaps whose midpoints are
exactly the elements of $A$. Adjoining an element $k\in A$ to the mark set splits the
$2$-gap centred at $k$ into two $1$'s; distinct elements of $U$ split distinct gaps.
Hence the content multiset of $\alpha(A^c\cup U)$ is $(2^{\,m-|U|},1^{\,e+2|U|})$, and
\[
  K_{\la,\alpha(A^c\cup U)}=\inner{s_\la,\,h_2^{\,m-|U|}h_1^{\,e+2|U|}}.
\]
Substituting into \eqref{eq:U-sum} and pulling the pairing outside,
\[
  G_\la(-q)=\Big\langle s_\la,\ \sum_{U\subseteq A}
    h_2^{\,m-|U|}h_1^{\,e+2|U|}\prod_{k\in U}\big(-q^{2k-1}\big)
    \prod_{k\in A\setminus U}\big(1+q^{2k-1}\big)\Big\rangle .
\]
The sum over $U$ factorises over $A$: for each $k\in A$ the choice $k\in U$ contributes
$h_1^{2}\cdot(-q^{2k-1})$ and the choice $k\notin U$ contributes
$h_2\cdot(1+q^{2k-1})$, the common factor $h_1^{e}$ being independent of $U$. Therefore
\[
  G_\la(-q)=\Big\langle s_\la,\ h_1^{e}\prod_{k\in A}
     \big(h_2(1+q^{2k-1})-h_1^2\,q^{2k-1}\big)\Big\rangle .
\]
Finally $h_2(1+q^{2k-1})-h_1^2q^{2k-1}=h_2+(h_2-h_1^2)q^{2k-1}=h_2-e_2\,q^{2k-1}$, using
$h_1^2=h_2+e_2$. This is the claim.
\end{proof}

\begin{remark}
At $q=1$ each factor is $h_2-e_2=p_2$, so
$G_\la(-1)=\inner{s_\la,h_1^e p_2^m}=\inner{s_\la,p_1^e p_2^m}=\chi^\la(2^m1^e)
=\chi^\la(w_0)$, the longest-element character: Proposition~\ref{prop:master}
contains and re-proves the previously established $q=-1$ fiber.
\end{remark}

\section{The vanishing law}

Two standard facts about the $2$-core are used; both are classical
(James--Kerber, \emph{The Representation Theory of the Symmetric Group}, \S2.7,
the theory of the $2$-quotient and Littlewood's formula).

\begin{lemma}[$2$-core obstruction]\label{lem:vanish}
If $b<|2\text{-}\core(\la)|$ then $\chi^\la(2^a1^b)=0$ (where $2a+b=n$).
\end{lemma}

\begin{proof}[Justification]
By Murnaghan--Nakayama, $\chi^\la(2^a1^b)=\sum_\mu(\pm)\,f^\mu$ summed over partitions
$\mu$ of size $b$ obtained from $\la$ by removing $a$ dominoes. Domino removal
preserves the $2$-core, so any such $\mu$ has the same $2$-core as $\la$ and hence
$|\mu|=b\ge|2\text{-}\core(\la)|$. If $b<|2\text{-}\core(\la)|$ there are no such
$\mu$ and the sum is empty.
\end{proof}

\begin{lemma}[$2$-quotient value]\label{lem:quotient}
Let $c_0=|2\text{-}\core(\la)|$ and $w=(n-c_0)/2$. Then
\[
  \chi^\la\!\big(2^{w}1^{c_0}\big)
  =\varepsilon_\la\; f^{\la^{(0)}}f^{\la^{(1)}}\binom{w}{|\la^{(0)}|}\,
   f^{\,2\text{-}\core(\la)},\qquad \varepsilon_\la\in\{\pm1\}.
\]
In particular this value is nonzero.
\end{lemma}

Define, for $0\le j\le m$,
\begin{equation}\label{eq:Ej}
  E_j:=\inner{s_\la,\,(-e_2)^{j}\,p_1^{e}\,p_2^{\,m-j}} .
\end{equation}

\begin{lemma}\label{lem:Ej}
$E_j=0$ for all $j<c(\la)$, and
$E_{c}=\big(-\tfrac12\big)^{c}\,\chi^\la\!\big(2^{m-c}1^{c_0}\big)\neq0$.
\end{lemma}

\begin{proof}
Using $e_2=\tfrac12(p_1^2-p_2)$,
\[
  (-e_2)^j p_1^e p_2^{m-j}
  =\Big(-\tfrac12\Big)^{j}\sum_{a=0}^{j}\binom ja(-1)^{j-a}\,p_1^{2a+e}\,p_2^{\,m-a},
\]
so that, pairing with $s_\la$,
\[
  E_j=\Big(-\tfrac12\Big)^{j}\sum_{a=0}^{j}\binom ja(-1)^{j-a}\,
       \chi^\la\!\big(2^{\,m-a}1^{\,2a+e}\big).
\]
The class $2^{m-a}1^{2a+e}$ has $2a+e$ fixed points. By Lemma~\ref{lem:vanish} its
character vanishes whenever $2a+e<c_0=2c+e$, i.e.\ whenever $a<c$. For $j<c$ every
index $a\le j$ satisfies $a<c$, so all terms vanish and $E_j=0$. For $j=c$ only the
top term $a=c$ survives, giving
$E_{c}=(-\tfrac12)^{c}\chi^\la(2^{m-c}1^{c_0})$, which is nonzero by
Lemma~\ref{lem:quotient}.
\end{proof}

\begin{proof}[Proof of Theorem~\ref{thm:main}]
Expand Proposition~\ref{prop:master} about $q=1$. Each factor is
\[
  h_2-e_2\,q^{2k-1}=p_2-e_2\,g_k,\qquad g_k:=q^{2k-1}-1=\sum_{i\ge1}\binom{2k-1}{i}(q-1)^i ,
\]
where we used $h_2-e_2=p_2$; note $g_k$ has no constant term, $g_k=(2k-1)(q-1)+O((q-1)^2)$.
Expanding the product over $A$,
\[
  h_1^e\prod_{k\in A}(p_2-e_2 g_k)
  =h_1^e\sum_{S\subseteq A}(-e_2)^{|S|}p_2^{\,m-|S|}\prod_{k\in S}g_k .
\]
Since $\prod_{k\in S}g_k$ is divisible by $(q-1)^{|S|}$, the coefficient of $(q-1)^r$
receives contributions only from $|S|=j\le r$, and pairing with $s_\la$,
\begin{equation}\label{eq:coeff}
  [\,(q-1)^r\,]\,G_\la(-q)
  =\sum_{j=1}^{r} E_j\!\!\sum_{\substack{S\subseteq A\\ |S|=j}}\!
    [\,(q-1)^r\,]\!\prod_{k\in S}g_k ,
\end{equation}
using \eqref{eq:Ej} and $h_1^e=p_1^e$, $\inner{s_\la,(-e_2)^jp_1^ep_2^{m-j}}=E_j$.
(The $j=0$ term contributes only to $r=0$.)

By Lemma~\ref{lem:Ej}, $E_j=0$ for $1\le j\le r$ whenever $r<c$; hence
$[(q-1)^r]G_\la(-q)=0$ for all $r<c$. For $r=c$ only $j=c$ survives in
\eqref{eq:coeff}, and for $|S|=c$ one has
$[(q-1)^c]\prod_{k\in S}g_k=\prod_{k\in S}(2k-1)$ (the lowest-order term of each
$g_k$). Summing over $S$,
\[
  [\,(q-1)^c\,]\,G_\la(-q)=E_{c}\,e_{c}\!\big(\{2k-1:k\in A\}\big)\neq 0,
\]
by Lemma~\ref{lem:Ej} and positivity of $e_c$. Since
$G_\la(q)=\sum_k a_k(q+1)^k$ gives
$G_\la(-q)=\sum_k a_k(-1)^k(q-1)^k$, we have
$[(q-1)^k]G_\la(-q)=(-1)^k a_k$. Therefore $a_k=0$ for $k<c$ and
\[
  a_{c}=(-1)^{c}E_{c}\,e_{c}\!\big(\{2k-1:k\in A\}\big)
       =(-1)^{c}e_{c}\!\big(\{2k-1:k\in A\}\big)\Big(-\tfrac12\Big)^{c}\chi^\la\!\big(2^{m-c}1^{c_0}\big)\neq0.
\]
Thus $\ord_{q=-1}G_\la=c(\la)$, and the displayed value is the leading coefficient.
Substituting Lemma~\ref{lem:quotient} for $\chi^\la(2^{m-c}1^{c_0})$ and absorbing the
sign $\varepsilon_\la$ together with the $(-1)$'s into $\pm$ yields the
$2$-quotient form, with
$S(n,c)=\pm\,e_c(\{2k-1:k\in A\})\,2^{-c}\,f^{\,2\text{-}\core(\la)}$ depending only on
$(n,c)$. The normalisations $S(n,0)=1$ (then $c=0$, $e_0=1$, core trivial) and
$S(n,1)=-\binom n2$ (then $e_1(\{2k-1:k\in A\})=\binom n2$ for $n$ odd and
$f^{(2,1)}=2$) are immediate.
\end{proof}

\section{The $c=1$ stratum, worked out}

Let $\la$ have $2$-core $(2,1)$, so $n$ is odd, $c_0=3$, $c=1$, $m=(n-1)/2$ and
$w=m-1=(n-3)/2$ dominoes. Here $A=\{2,4,\dots,n-1\}$ and
$e_1(\{2k-1:k\in A\})=\sum_{k\in A}(2k-1)=\binom n2$. Since the $2$-core is nonempty,
$\chi^\la(w_0)=G_\la(-1)=0$ (Lemma~\ref{lem:vanish}), so the order is at least $1$;
Theorem~\ref{thm:main} gives order exactly $1$ with
\[
  a_1=-E_1\,\binom n2=\tfrac12\,\binom n2\,\chi^\la\!\big(2^{m-1}1^{3}\big)
     =(-1)^{n(\la)+1} f^{\la^{(0)}}f^{\la^{(1)}}\binom{w}{|\la^{(0)}|}\binom n2,
\]
the last step by Lemma~\ref{lem:quotient} (which here gives
$\chi^\la(2^{m-1}1^3)=(-1)^{n(\la)+1}\cdot2\cdot
f^{\la^{(0)}}f^{\la^{(1)}}\binom{w}{|\la^{(0)}|}$, with $f^{(2,1)}=2$). For the
smallest case $\la=(2,1)$ ($n=3$, $w=0$): $G_{(2,1)}(q)=q^3+1=(q+1)(q^2-q+1)$, so
$a_1=3$, matching $(-1)^{1+1}\cdot1\cdot1\cdot\binom00\cdot\binom32=3$.

\section{Computational verification}

Every claim was checked by independent computation (exact integer / rational
arithmetic), with no exceptions, for all partitions in the indicated range.
\begin{itemize}
\item \textbf{Master identity (Proposition~\ref{prop:master}).} The polynomial
identity $G_\la(-q)=\inner{s_\la,h_1^e\prod_{k\in A}(h_2-e_2q^{2k-1})}$, with the
right side expanded into power sums and evaluated by Murnaghan--Nakayama, agrees with
direct $\SYT$ enumeration of $G_\la$ for all $\la\vdash n$, $2\le n\le 9$.
\item \textbf{Order law and leading coefficient.} For all $\la\vdash n$,
$4\le n\le 12$ ($265$ shapes), $\ord_{q=-1}G_\la=c(\la)$, and the leading Taylor
coefficient $a_c$ equals
$(-1)^c e_c(\{2k-1:k\in A\})(-\tfrac12)^c\chi^\la(2^{m-c}1^{c_0})$.
\item \textbf{$E_j$ vanishing (Lemma~\ref{lem:Ej}) and the $2$-core obstruction
(Lemma~\ref{lem:vanish}).} $\chi^\la(2^a1^b)=0$ for all $b<|2\text{-}\core(\la)|$, and
$E_j=0$ for $j<c$, for all $\la\vdash n\le 12$.
\item \textbf{$S(n,c)$.} The reduced leading coefficient
$a_c/\big((-1)^{n(\la)}f^{\la^{(0)}}f^{\la^{(1)}}\binom{w}{|\la^{(0)}|}\big)$ is
single-valued for each $(n,c)$ with $n\le12$, and equals
$\pm e_c(\{2k-1:k\in A\})2^{-c}f^{\,2\text{-}\core}$ (e.g.\ $S(n,1)=-\binom n2$,
$S(6,3)=90$, $S(8,3)=1624$, $S(10,5)=238680$, $S(12,5)=7490736$).
\end{itemize}

\paragraph{Scripts.}
\texttt{\~{}/projects/scratch/prove-2026-06-01-graded-2core/} (notably
\texttt{verify\_master.py}, \texttt{verify\_general.py}, \texttt{verify\_full.py},
\texttt{referee\_check.py}).

\end{document}
